\subsection{Protenix output remapping with the Factor head frozen}
\label{app:followup_map}

The OpenFold compensation study above trained the residual generator and
output map jointly. To test whether remapping alone can improve a
trained head, we started from the 12 Protenix Factor fits in
Appendix~\ref{app:followup_seeds}: one Native and three Rotated fits
for each of three training seeds. We forked each parent into three
512-update branches: map only (frozen head), head only (fixed identity
map), and joint training. The branches used the same continuation
target order and paired training noise; the two branches that updated
the head inherited its optimizer state.

Each branch with a trainable map learned its own identity-initialized
$C\in SO(128)$ via the exponential of a full skew-symmetric matrix.
The map is shared across targets, pairs, and recycles within that model.
It is applied after the fixed rotation and before adding the query
baseline, and preserves each pair's residual norm without a channel-mean
constraint. We evaluated the new branches on the previously
evaluated Fresh192 panel with the same three inference seeds as the
parent predictions. For the map-only branch, let $B_N$ and $B_R$ be
its paired gains over the corresponding parents, with the latter
averaged over the three rotations. The primary contrast is $B_R-B_N$.

\begin{table}[!htbp]
\centering\small
\setlength{\tabcolsep}{4pt}
\caption{Map-only gains in C$\alpha$ pair-lDDT with the trained
Factor head frozen. The 95\% intervals use 20,000 paired bootstrap
draws, resampling targets within the fixed Fresh192 length strata.
The exploratory column additionally resamples the three training
seeds; inference seeds and rotations stay fixed.}
\label{tab:followup_map}
\begin{tabular}{@{}lrrr@{}}
\toprule
Contrast & Estimate & Target CI & Target and seed CI \\
\midrule
\rebuttalFrozenMapRows
\bottomrule
\end{tabular}
\end{table}

\begin{table}[!htbp]
\centering\small
\caption{C$\alpha$ pair-lDDT for the same Protenix parents and continuation
branches. Contrast intervals resample targets only and are unadjusted.}
\label{tab:followup_map_controls}
\begin{tabular}{@{}lrr@{}}
\toprule
Branch or contrast & Native & Rotated \\
\midrule
\rebuttalBranchRows
\bottomrule
\end{tabular}
\end{table}

The map improved the rotated head with the residual generator frozen.
Native Factor also improved. The interval for the extra Rotated gain
barely excluded zero when resampling targets alone and included zero
when resampling training seeds as well (Table~\ref{tab:followup_map}).
This intervention shows that remapping alone can improve predictions
under this recipe. It leaves the cause of the original head-by-rotation
interaction unresolved.
